\documentclass[10pt]{article}
\usepackage[margin=1in]{geometry}
\usepackage{amsmath,amssymb}
\title{Disproof of Conjecture 1665: the proposed tight burning family}
\author{gaochengzhecpu (AI-assisted submission)}
\date{October 3, 2026}
\usepackage{url}
% Compact consistent title for a short mathematical note.
\makeatletter
\renewcommand{\maketitle}{\begin{center}
{\Large\bfseries\@title\par}\vspace{0.6em}
{\small\@author\par\@date}
\end{center}\vspace{0.5em}}
\makeatother
\setlength{\emergencystretch}{3em}
\begin{document}
\maketitle

\paragraph{Claim being disproved.}
The conjecture states, among other claims, that Cartesian products of the
Petersen graph with paths have burning number
$\lceil\sqrt{n}\rceil+1$, where $n$ is the order of the product.
We disprove this universal tight-family statement. We make no assertion
about the separate proposed upper bound for arbitrary graphs.

\paragraph{Definitions and graph.}
Let $P$ be the Petersen graph on vertices $0,\ldots,9$. Its edges are
the outer cycle $i(i+1)$ for $i\in\mathbb Z/5\mathbb Z$, the inner
cycle $(5+i)(5+i+2)$, and the five spokes $i(5+i)$.
Let $P_2$ be the path with two vertices $0,1$. Set $G=P\mathbin\square P_2$.
Write the product vertex $(u,j)$ as the integer $u+10j$.
Thus $G$ is a simple graph of order $20$.

At each burning round, choose a vertex unburned at the beginning of the
round, ignite it, and spread the older fires to all adjacent vertices.
For the burned set $B$ and source $s$, the update is
\[
 B\longmapsto B\cup\{s\}\cup\{v: (\exists u\in B)\ uv\in E(G)\}.
\]
The burning number is the least number of rounds in which every vertex
can be burned.

\paragraph{A four-round certificate.}
Choose the successive sources $0,1,2,12$. Direct application of the
update rule gives
\begin{align*}
 B_0&=\varnothing,\\
 B_1&=\{0\},\\
 B_2&=\{0,1,4,5,10\},\\
 B_3&=\{0,1,2,3,4,5,6,7,8,9,10,11,14,15\},\\
 B_4&=\{0,1,\ldots,19\}.
\end{align*}
Each source is outside the burned set at the beginning of its round:
$0\notin B_0$, $1\notin B_1$, $2\notin B_2$, and $12\notin B_3$.
Consequently $b(G)\leq4$.
Equivalently, the Petersen graph has diameter two and its product with
$P_2$ has diameter three, so one initial fire already reaches every
vertex by the fourth round.

On the other hand, $4^2<20\leq5^2$, so
$\lceil\sqrt{20}\rceil+1=6$. Therefore
\[
 b(P\mathbin\square P_2)\leq4<6
 =\lceil\sqrt{|V(P\mathbin\square P_2)|}\rceil+1.
\]
This contradicts the claimed equality and disproves Conjecture 1665.

\paragraph{Lean verification and semantic correspondence.}
The self-contained file \texttt{lean/Main.lean} imports only \texttt{Std}.
The predicate \texttt{petPath m} is the Cartesian product with the
$m$-vertex path. Theorems \texttt{product\_encoding} and
\path{product_encoding_bijective} verify that the 20 labels used
here represent exactly $P\square P_2$, and \texttt{graph\_simple}
checks symmetry and absence of loops.
The definitions \texttt{advance}, \texttt{run}, and \texttt{legal}
implement precisely the update rule above, including the condition on
new sources. Theorems \texttt{legal\_sources} and \texttt{all\_burned}
check the complete finite process, not merely its length.
The predicate \texttt{BurningNumberIs} asserts successful burning and
minimality. Finally, \path{not_tight_family_claim} refutes the
universally quantified equality for all products with nonempty paths.
Its reported axiom dependency is only \texttt{propext}. No
\texttt{sorry}, \texttt{admit}, \texttt{native\_decide}, or additional
axiom is used.

\paragraph{Reproduction.}
Run \texttt{lean Main.lean} with Lean 4.19.0. The command exits
successfully and prints the axiom dependencies of the principal results.
All finite certificates use kernel-checked \texttt{decide}.
\end{document}
